\chapter{The Media System}
\label{ch:media-system}

\section{A correction worth stating up front}

This book's opening chapter quoted CNA's own README verbatim: the \cnans{Media} namespace is
``25/25 present, but 14 are shells.'' That figure is real, but it is also, at the time of
writing, \textbf{stale} --- confirmed by directly reading the namespace's own current
tracking record. A dedicated, from-scratch effort ran sixteen phases and 231 tasks against
this namespace, closing 224 of them, backed by fifteen independent adversarial external
review rounds. All 24 \cnans{Media} types are now real implementations backed by genuine file
parsing --- Vorbis and ID3v2 tag parsing, \texttt{.m3u} / \texttt{.m3u8} playlist parsing,
FFmpeg-backed video, SDL3\_mixer-backed audio --- not shells. The current media module contains
285 test cases. The formerly quoted 5{,}471/5{,}467 figure was a whole-suite snapshot with a
Devices-specific hardware-skip note attached to the wrong namespace. It should not be replaced
by another context-free total: Chapter~\ref{ch:counting-discipline} derives separate source,
configured-build, registration, and runtime populations and explains why none is the universal
``current test count.''
This is exactly the kind of documentation drift this book's first chapter warned readers to
watch for: a specific, dated, checkable claim that this book preserves as evidence of the
project's own earlier state, corrected here against the namespace's own current record rather
than silently repeated as still current.

Audio and Media form an explicit static-library cycle: the audio-side
\cnaclass{FrameworkDispatcher} pumps \cnaclass{MediaPlayer}, while Media routes song playback
through the audio mixer. CMake records both edges and resolves the archive cycle by repeating
the libraries. This is genuine XNA-semantic coupling, not an accidental include dependency;
breaking it requires moving ownership of the shared dispatcher/pump contract, not deleting one
link line.

\section{Song, and the one gap even FNA has}

\cnaclass{Song} is a real, file-backed asset, constructed either from an explicit path or
discovered through a library scan (\S\ref{sec:medialibrary}). \texttt{Album} / \texttt{Artist}/%
\texttt{Genre} back-references are non-owning and remain null unless the song was actually
discovered by a scan rather than constructed directly. \texttt{IsRated} / \texttt{Rating}
(on XNA's 0--10 scale) are populated from real ID3v2 \texttt{POPM} or Vorbis \texttt{RATING}
tags, with \texttt{IsRated} deliberately reporting false for an explicit rating of zero,
since both tag formats reserve that value to mean ``unrated'' rather than ``rated zero.''
\texttt{IsProtected} is always false --- correctly, since a DRM-wrapped file is not something
a local library scan could index in the first place.

One finding here is worth its own paragraph, because of what it reveals about verifying
against the wrong reference: FNA's own \texttt{Song} class omits \texttt{Album},
\texttt{Artist}, \texttt{Genre}, and \texttt{ToString()} relative to the real XNA~4.0
reference assemblies --- a gap invisible across eight separate audit rounds, because every
one of them diffed against FNA's source rather than against Microsoft's own original
reference documentation directly. It was only found by finally diffing against the real
XNA~4.0 API surface itself (the same \texttt{xna4-spec} resource covered in
Chapter~\ref{ch:xna4-spec-auditing}). This is the sharpest illustration in this entire book of
why this book's first chapter insisted on distinguishing ``matches FNA'' from ``matches real
XNA'' as two different claims: here, they were not the same claim at all, and treating FNA as
an infallible stand-in for the real spec cost eight rounds of review before the gap surfaced.

\subsection{Full property reference, and two deliberate deviations}

Beyond the library back-references already covered, \cnaclass{Song} exposes
\texttt{Duration} (settable, since \cnaclass{MediaPlayer} rewrites it once the real decoded
duration is known --- see below), \texttt{PlayCount} (also settable, and reset by
\cnaclass{MediaPlayer}\texttt{::Stop()}), and \texttt{Handle}, a CNAEXT extension exposing the
backend file path or handle a worked example needing direct file access can read without
going through \cnaclass{MediaLibrary}. \texttt{Equals} compares by handle, and
\texttt{GetHashCode} is deliberately \emph{content}-based --- a hash of the resolved file
handle --- in explicit, documented contrast to FNA's own identity-based
\texttt{base.GetHashCode()}. That FNA choice means two FNA \cnaclass{Song}s that are
\texttt{Equals}-equal (same handle) can still report different hash codes, violating the usual
Equals/GetHashCode contract; CNA does not replicate that inconsistency, on the reasoning that a
known-broken FNA behavior is not owed parity the way a genuine XNA behavior is.

\begin{lstlisting}
// Worked example: constructing a standalone Song outside any MediaLibrary
// scan (Album/Artist/Genre all stay nullptr, matching a song with no
// library context), then handing it directly to MediaPlayer. FromUri()
// returns an owning raw pointer; keep that ownership explicit. Play()
// clones the Song into its own queue rather than adopting this pointer.
std::unique_ptr<Song> introTrack(
    Song::FromUri("Intro Theme", "Content/Audio/intro.ogg"));
MediaPlayer::Play(introTrack.get());

// Duration is whatever the caller supplied (zero, if none) until playback
// actually starts -- MediaPlayer overwrites it with the real decoded
// duration read back from the audio backend once the file is opened.
\end{lstlisting}

\subsection{The resolver advertises two undecodable extensions}

\cnaclass{MediaLibrary} correctly avoids indexing AAC and WMA, but the loose-file
\texttt{SongTypeReader} separately advertises \texttt{.aac} and \texttt{.wma} alongside MP3,
Ogg, WAV, FLAC, and Opus. CNA's vendored SDL3\_mixer build does have an Opus decoder, and the
library scanner deliberately includes \texttt{.opus}; it has no AAC or WMA decoder.
Loose-file resolution can therefore select an AAC or WMA path that construction can represent
but the playback backend cannot decode. This is a real inconsistency between extension
discovery and the playback dependency, not a promise that every listed extension works.

\section{The library hierarchy}
\label{sec:medialibrary}

\cnaclass{SongCollection}, \cnaclass{Album} / \texttt{AlbumCollection},
\cnaclass{Artist} / \texttt{ArtistCollection}, and \cnaclass{Genre} / \texttt{GenreCollection}
form a real, populated hierarchy, built by \cnaclass{MediaLibrary} through a genuine
filesystem scan of the host's music and pictures roots. \texttt{Album::GetAlbumArt()}/%
\texttt{GetThumbnail()} return a real \texttt{System::IO::Stream}, resolved either by
filename convention (\texttt{cover.jpg}, \texttt{folder.jpg}) or extracted on demand from an
embedded ID3v2 \texttt{APIC} or FLAC \texttt{PICTURE} frame --- extraction directly from an
embedded frame at scan time, rather than on demand, is a documented, deliberate future
increment, not a bug. \cnaclass{Picture} / \texttt{PictureAlbum} and their collections form a
parallel, real hierarchical tree for image libraries; \texttt{MediaLibrary::SavePicture()}
lazily creates a real ``Saved Pictures'' node the first time it is needed.

\subsection{How the scan actually reads tags}

The host's Music and Pictures roots themselves are resolved through SDL's own
\texttt{SDL\_GetUserFolder()}, giving each platform its real, conventional folder rather than a
hardcoded path guess --- with a CNAEXT static override hook specifically so the test suite does
not depend on whatever music a given CI machine's real user account happens to have. Tag
reading is a genuine, from-scratch, dependency-free parser, not a third-party library
wrapper: real Ogg Vorbis-comment parsing, real ID3v2.3/2.4 text-frame parsing (\texttt{TIT2}/%
\texttt{TPE1} / \texttt{TALB} / \texttt{TCON} / \texttt{TRCK}, including correct handling of all four
ID3v2 text-encoding bytes --- Latin-1, UTF-16 with a byte-order mark, UTF-16BE, and UTF-8), and
--- beyond what this chapter's introduction above already covers --- real native-FLAC
\texttt{VORBIS\_COMMENT} metadata-block parsing and real Ogg-Opus \texttt{OpusTags} comment-header
parsing as well, so both of those container formats get real tags too, not just Ogg Vorbis and
MP3. A file with no readable tags at all --- or an unsupported format --- falls back to
deriving \texttt{Title} / \texttt{Album} / \texttt{Artist} from the file's own path (filename,
parent directory, and grandparent directory respectively), so a scan never simply drops an
otherwise-playable file for lacking metadata.

\texttt{TRCK} deserves its own note, because the fix here is a different shape of bug than
anything else in this section --- not a missing parser, but a parsed value silently discarded
one layer up. The tag parser itself always genuinely read and stored a real track number; the
bug was that \cnaclass{MediaLibrary}'s own \cnaclass{Song}-construction call only ever passed
through \texttt{(path, title)} or \texttt{(path, title, durationMs)}, with no parameter slot
for the track number at all, so \texttt{Song::TrackNumber} stayed hardcoded at~0 regardless of
what the file actually tagged. The real, mutation-verified fixture that caught this uses four
distinct, independently-confirmed track numbers on purpose --- \emph{Sunrise}=1,
\emph{Twilight}=1, \emph{Daybreak}=2, \emph{{\'E}toile}=3 --- specifically so an
all-zero-or-all-one stub implementation cannot accidentally pass by coincidence the way a
single-fixture test might. This is worth remembering as a general testing lesson beyond this
one property: a bug that silently drops a correctly-parsed value one layer up in the call chain
is invisible to any test that only checks the parser in isolation, and a fixture whose expected
values are not all identical is what actually forces a real, non-stub implementation to exist.

The user-rating conversion is a real, if underspecified, corner: ID3v2's \texttt{POPM} frame
uses a 0--255 byte, while the Vorbis \texttt{RATING} comment has no agreed standard at all ---
different taggers write 0--100, 0--5, or 0--10 into it in practice. Neither real XNA nor FNA
defines a conversion for either case, since neither one reads tags from arbitrary files at all;
the scale mapping CNA applies to reach the 0--10 range \cnaclass{Song::Rating} exposes is
therefore a CNA decision, recorded as such in the namespace's own tracking documentation rather
than presented as a spec-mandated conversion. Embedded cover-art extraction (ID3v2
\texttt{APIC}, or the FLAC \texttt{METADATA\_BLOCK\_PICTURE} block type~6) prefers the
front-cover picture type when a file embeds more than one image, falling back to whichever
image comes first otherwise.

Playlist scanning tolerates real-world messiness the same way: a \texttt{.m3u} / \texttt{.m3u8}
entry pointing at a file that no longer exists is silently skipped, not treated as a fatal
parse error, matching FNA's own general ``missing content degrades gracefully'' posture; and
both extensions are read identically, as raw UTF-8 bytes, with no separate legacy-encoding path
for plain \texttt{.m3u} --- verified directly against a non-ASCII fixture file name to confirm
the round-trip actually holds, not merely assumed from the two formats sharing a parser.

\section{Playback: MediaPlayer, MediaQueue, Playlist}

\cnaclass{MediaPlayer} is the static playback controller familiar from real XNA:
\texttt{Play(Song*)} / \texttt{Play(SongCollection)}, \texttt{Pause} / \texttt{Resume}/%
\texttt{Stop}, a real \cnaclass{MediaQueue} with shuffle and repeat, and
\texttt{GetVisualizationData()}. Playback position tracking is real, timer-based state, with
an explicit fallback song-end detector for cases where no native ``track finished'' signal is
available from the underlying decoder. \cnaclass{Playlist} / \texttt{PlaylistCollection}
support genuine, read-only \texttt{.m3u} / \texttt{.m3u8} parsing --- there is no write API, but
that matches real XNA's own \cnaclass{Playlist} being read-only too, not a CNA-specific
limitation.

\subsection{MediaQueue's ownership model, and why Play() clones its songs}

\cnaclass{MediaQueue}\texttt{::Add(Song*)} takes ownership of the pointer it's given --- reading
\texttt{MediaQueue.cpp} directly confirms it wraps every added song in a
\texttt{std::unique\_ptr} internally, and \texttt{Clear()} (called at the start of every
\texttt{Play()}) destroys whatever the queue currently owns. This creates a real hazard
\cnaclass{MediaPlayer} has to defend against: \cnaclass{SongCollection}, the type
\texttt{Play(const SongCollection\&)} accepts, explicitly does \emph{not} own the songs it
holds --- it is a thin, non-owning view over songs that belong to a \cnaclass{MediaLibrary} (or
to the caller). Handing those same non-owning pointers directly to an owning queue would mean
the very next \texttt{Play()} call's \texttt{Clear()} deletes songs the library still expects
to be alive. \cnaclass{MediaPlayer} avoids this by never storing the caller's own
\cnaclass{Song} pointers at all: its private \texttt{LoadSong()} helper constructs a fresh,
independent \texttt{new Song(song->getHandle(), song->getNameProperty())} clone for every song
it enqueues, and it is that clone the queue owns and eventually destroys --- the original,
externally-owned \cnaclass{Song} is never touched.

A second, easy-to-miss behavior: \texttt{MediaPlayer::Stop()} does not just stop playback ---
it walks the entire current queue and resets every song's \texttt{PlayCount} back to zero, not
only the song that was actually playing. And when no native ``track finished'' signal is
available from the underlying decoder (a build without \texttt{SOUND\_ENABLED}), the
\texttt{Update()} fallback branch calls \texttt{DetectSongEndedByElapsedTime} exactly once per
frame, comparing elapsed playback time against the song's own \texttt{Duration} and only
reporting ``ended'' once that duration is genuinely known (greater than zero) --- a
\cnaclass{Song} with an unset, zero \texttt{Duration} simply never auto-advances this way, on
the reasoning that a false trigger at time zero would be worse than never advancing at all. A
normal, sound-enabled build instead relies on a real \texttt{MIX\_SetTrackStoppedCallback} fired
from the audio thread, so the elapsed-time fallback exists purely for the no-audio build
configuration, not as the primary end-of-track mechanism.

\begin{lstlisting}
// Worked example: MediaPlayer::Stop() resets play counts for the WHOLE
// queue, not just the song that was playing -- easy to miss if you only
// read the method name.
SongCollection* songs = library.getSongsProperty();
MediaPlayer::Play((*songs)[0]);
// ... song plays through, MediaPlayer::MoveNext() advances a few times ...
MediaPlayer::Stop();
// Every song MediaPlayer cloned into its queue now reports PlayCount() == 0
// again, even ones that had already finished playing earlier in the queue.
\end{lstlisting}

\texttt{MoveNext()} / \texttt{MovePrevious()} both stop the current track first, then pick the
next index: shuffled mode draws a uniform random index over the whole queue (so the same song
can, correctly, play again back-to-back by chance --- this matches a simple ``pick any index''
shuffle, not a shuffle-bag that guarantees no repeats), while linear mode clamps the index
within bounds unless repeat is on, in which case reaching the last song wraps back to index
zero instead of clamping there. Streamed playback itself does not fully predecode a song file
up front (the backend is asked to stream it), which matters for a long music track's memory
footprint; and \cnaclass{Song}\texttt{::Duration} is deliberately overwritten with the real
duration read back from the decoded audio asset once playback actually starts, regardless of
whatever value --- including zero --- the \cnaclass{Song} was constructed with.

\subsection{Volume and mute are independent state, and events fire late on purpose}

\cnaclass{MediaPlayer}\texttt{::Volume} and \texttt{IsMuted} read as if muting might simply
force the volume to zero, but reading \texttt{setVolumeProperty()} /
\texttt{setIsMutedProperty()} directly in \texttt{MediaPlayer.cpp} shows they are genuinely
independent state, both funneled through the same private \texttt{ApplyMusicVolume(volume,
muted)} helper:

\begin{lstlisting}
void MediaPlayer::setVolumeProperty(float value)
{
    volume_ = MathHelper::Clamp(value, 0.0f, 1.0f);
    ApplyMusicVolume(volume_, isMuted_);   // muted flag re-applied, volume_ untouched
}

void MediaPlayer::setIsMutedProperty(bool value)
{
    isMuted_ = value;
    ApplyMusicVolume(volume_, isMuted_);   // volume_ re-applied, isMuted_ untouched
}

// Inside ApplyMusicVolume: MIX_SetTrackGain(track, muted ? 0.0f : vol);
\end{lstlisting}

\texttt{volume\_} is never overwritten by muting, and \texttt{isMuted\_} is never cleared by
changing the volume --- only the actual SDL3\_mixer gain call, \texttt{MIX\_SetTrackGain()},
substitutes \texttt{0.0f} for the real value while muted. A game that calls
\texttt{setVolumeProperty(0.3f)} while already muted, then later calls
\texttt{setIsMutedProperty(false)}, hears \texttt{0.3f} immediately --- the volume set while
silent was never lost, and unmuting does not require the caller to remember and re-apply it.

A second, easy-to-miss design decision: \texttt{MediaStateChanged} and \texttt{ActiveSongChanged}
do not fire the instant \texttt{setStateProperty()} or a queue-advance actually changes
something. Both setters only flip a static flag on \cnaclass{FrameworkDispatcher}
(\texttt{FrameworkDispatcher::MediaStateChanged = true;}); the events themselves are raised
later, from \texttt{FrameworkDispatcher::Update()}, which checks the flag, calls
\texttt{MediaPlayer::OnMediaStateChanged()} if it is set, and clears it again. This is not a
CNA-specific quirk --- real XNA's own \cnaclass{FrameworkDispatcher}\texttt{.Update()} is a
documented API a game must call once per frame, and CNA's \cnaclass{Game} base class already
calls it automatically as part of its own update cycle, so an ordinary game observes no
difference from immediate dispatch. The real, checkable consequence is that changing
\texttt{MediaPlayer} state and then reading the event's own fired-or-not status \emph{before}
the next \texttt{FrameworkDispatcher::Update()} call would see it not yet fired --- exactly
what the real regression test for this confirms, by driving the check through the actual
dispatcher call chain rather than invoking \texttt{OnMediaStateChanged()} directly:

\begin{lstlisting}
Song song(kFixtureA, "A");
MediaPlayer::Play(&song);  // flips both flags: new active song + Stopped -> Playing

Microsoft::Xna::Framework::FrameworkDispatcher::Update();  // events actually fire HERE

// Only after Update() do the ActiveSongChanged/MediaStateChanged subscribers see anything.
\end{lstlisting}

\section{Video and VideoPlayer}

\cnaclass{Video} / \texttt{VideoPlayer} are real, backed by genuine FFmpeg decoding, CNA's own
graphics-renderer frame rendering, and SDL3 audio streaming, on Linux and macOS.
\texttt{VideoPlayer::Play()} deliberately throws \texttt{InvalidOperationException} if a
video's own declared width, height, or frame-rate metadata does not match what the decoded
file itself actually reports --- a real, checked consistency guard rather than a silent
mismatch. \cnaclass{VideoSoundtrackType} (\texttt{Music} / \texttt{Dialog}/%
\texttt{MusicAndDialog}) is confirmed metadata-only in both FNA and CNA alike; neither
implements any actual audio-ducking or muting behavior tied to this value.

\subsection{A real disposed-state gap, closed with one deliberate, C\texttt{++}-specific exception}

\cnaclass{VideoPlayer} once had \texttt{getIsDisposedProperty()} but nothing that actually
consulted it: every public method touched the player's own state unconditionally after
\texttt{Dispose()}, whether or not anything was still safe to touch. The fix reproduces real
FNA's own \texttt{checkDisposed()} guard --- called at the top of eight public methods --- down
to its exact, faithfully-reproduced message: \texttt{throw
System::ObjectDisposedException("VideoPlayer")}, a hardcoded literal type-name string matching
FNA's own choice rather than anything reflection-derived, for message fidelity rather than mere
behavioral parity. One call site is a deliberate, reasoned exception rather than an oversight:
\texttt{Dispose()} itself does \emph{not} call this guard, even though real FNA's own
\texttt{Dispose()} does (throwing on a second call). The reason is specific to C\texttt{++}, not
a parity shortfall --- \texttt{\textasciitilde VideoPlayer()} unconditionally calls
\texttt{Dispose()}, and a destructor is implicitly \texttt{noexcept} in C\texttt{++}, so a second
explicit \texttt{Dispose()} call followed by ordinary destruction would throw from inside a
destructor and terminate the program outright, rather than merely surprising a caller the way
the equivalent double-\texttt{Dispose()} does in C\#. \texttt{Dispose()} is kept safely
idempotent instead, and every other public method keeps the real, faithful throw --- a small,
precise illustration of this book's own recurring theme: a documented, reasoned divergence from
FNA, not a silent one, and specifically motivated by a real language-level constraint C\# does
not share.

\subsection{Multi-track selection, and three real bugs an external review found}

Both \cnaclass{Video} and \cnaclass{VideoPlayer} expose \texttt{SetAudioTrackEXT(track)} /
\texttt{SetVideoTrackEXT(track)} --- an FNA extension beyond the original XNA~4.0 surface (note
the \texttt{EXT} suffix; not a CNA invention) for selecting among multiple audio or video
streams inside one container file. Switching either track live, mid-playback, is exactly the
scenario a real external code-review pass targeted, and it found three separate, genuine bugs
rather than a hypothetical concern:

\begin{itemize}
  \item A video playing with \emph{no} audio device available (or a video-only track selected)
  used to accumulate its entire decoded audio track in memory for the rest of playback ---
  potentially hundreds of megabytes for a long video --- because the old buffer-draining code
  only cleared its buffer inside the branch that checks for a live audio stream, never in the
  no-stream case, and \texttt{CloseDecoder()} did not clear it either. Worse, stale audio left
  over from a failed-device \texttt{Play()} could then be fed into a genuinely opened stream on
  a later, successful \texttt{Play()} call. The fix, \texttt{DrainAndFlushAudioBuffer()}, always
  clears the buffer after handing it off, whether or not a live stream exists to receive it.
  \item Switching video tracks mid-playback can change frame dimensions, but the existing frame
  texture used to silently keep its old size regardless --- a stale-size bug invisible until a
  file with differently-sized video tracks was actually tested with a live track switch.
  \texttt{ReconfigureVideoOutputForCurrentTrack()} now recreates the frame texture to match
  whichever track is actually active.
  \item Audio and video track reconfiguration used to be a single function always invoked
  together, so switching only the video track needlessly tore down and reopened the SDL audio
  stream too --- discarding whatever audio was already queued for playback --- and, symmetrically,
  an audio-only switch needlessly recreated the video texture. The fix splits the two into
  independent helpers, one per side: \texttt{ReconfigureAudioOutputForCurrentTrack()} handles
  the audio stream.

  \texttt{ReconfigureVideoOutputForCurrentTrack()} handles the frame texture. A track switch
  now only reconfigures whichever side actually changed.
\end{itemize}

\begin{lstlisting}
// Worked example: a container with two audio tracks (say, two dubbed
// languages) and one video track. Switching languages mid-playback no
// longer touches the video output at all, thanks to the split above.
VideoPlayer player;
player.Play(introVideo);
// ... later, the player picks a different language track ...
player.SetAudioTrackEXT(1);  // only the audio stream is torn down and reopened
\end{lstlisting}

\subsection{Visualization: a genuine, from-scratch pipeline, not a stub}

\texttt{MediaPlayer::GetVisualizationData()} is backed by a real, working pipeline end to end,
not a hardcoded or randomized placeholder. SDL3\_mixer's post-mix callback --- which runs on
the \emph{audio thread}, under a real-time constraint of no allocation, no locking, and no
exceptions --- feeds every mixed buffer into a lock-free, single-producer/single-consumer ring
buffer sized to hold more than one FFT window at a time. Reading directly from that ring
buffer's own header comment surfaces a genuinely subtle, verified finding: its sample storage
is \texttt{std::atomic<float>}, accessed with relaxed memory ordering, specifically because a
plain \texttt{float} written by the audio thread while the game thread reads it is a data race
and therefore undefined behavior in C++ regardless of how harmless the generated machine code
looks --- relaxed atomics make the same access well-defined at effectively zero runtime cost on
every mainstream platform, since a relaxed load or store of a naturally-aligned float compiles
to the exact same instruction a plain access would. What the design deliberately does
\emph{not} guard against is a reader observing a window whose samples span two different
callback batches --- a torn read in the logical sense, not the undefined-behavior sense --- and
that residual imprecision is accepted as a one-frame visual artifact in a visualizer, not
something worth a lock on the audio path.

The frequency-domain half comes from a genuine, dependency-free, from-scratch radix-2 FFT: a
512-real-sample transform (XNA exposes exactly 256 frequency bins for visualization, so a
textbook transform of twice that size is entirely sufficient, and pulling in a third-party FFT
library for it would be disproportionate), with a Hann window applied first to reduce spectral
leakage and a 2/512 magnitude-scaling factor chosen so that a full-scale sine wave maps
to approximately~1.0 in its own bin --- a CNA-chosen normalization recorded as such, since XNA
itself never documented its own scale for this API. Two further, real bugs an external review
found and fixed in the surrounding enable/disable logic are worth naming directly: an early
version set the ``visualization enabled'' flag \emph{before} even confirming a mixer could be
obtained, so a failed mixer acquisition left \texttt{IsVisualizationEnabled} reporting true with
no mixer and no callback actually installed; and a failed attempt to \emph{un}install the tap
used to leave the internal ``still running'' flag stuck false (matching the caller's request for
off) even though the callback might still genuinely be live, so a later retry never actually
tried again. \texttt{GetVisualizationData()} itself does not throw when visualization is off or
no data has been captured yet --- it fills both arrays with zero, matching
\cnaclass{VisualizationData}'s own zero-initialized construction rather than treating that state
as exceptional.

\section{The one genuinely open gap: Windows, Android, and Emscripten video}

The one area of this namespace still genuinely open, rather than merely undocumented, is
real FFmpeg-backed video on Windows, Android, and Emscripten --- excluded from those builds
entirely via a \texttt{CNA\_FFMPEG\_AVAILABLE} CMake gate. Using \texttt{Video}/%
\texttt{VideoPlayer} there today is a \emph{link error}, not a runtime stub throwing a
graceful exception --- a direct consequence of the project owner explicitly rejecting a
\texttt{NotSupportedException}-stub approach in favor of holding out for a real FFmpeg
integration on those platforms instead. This is not closeable from a Linux-only development
environment, and is recorded here as exactly that kind of gap: not silently missing, not
falsely claimed complete, but named and left for real cross-platform FFmpeg packaging work.

There is also a concrete native-Windows build defect in the two-layer gate. CMake disables the
video translation units for \texttt{WIN32}, MinGW, Emscripten, and Android, but
\texttt{ContentManager.cpp}'s matching include guard repeats \texttt{\_\_MINGW32\_\_} and never
checks native \texttt{\_WIN32}. An MSVC build can therefore retain the \cnaclass{Video} include
and dispatch reference while the implementation was excluded, producing a link failure. The
comment beside that guard also omits \texttt{WIN32} when describing the CMake rule. Whenever a
feature is gated in both build code and C++, test their truth tables rather than trusting two
similar-looking expressions.

\section{Deliberate, documented non-gaps}

A short list of choices that look like gaps but are not: \texttt{.m4a} / \texttt{.aac}/WMA
files are deliberately not indexed by \cnaclass{MediaLibrary} at all, because SDL3\_mixer
ships no AAC decoder --- indexing a song the player could never actually play would be worse
than not listing it. There is no \texttt{MediaLibrary::Refresh()}, because a library scan is
a deliberate one-time synchronous snapshot, not a live-updating view. And
\texttt{MediaPlayer::Update()}'s fallback branch for a build with sound disabled is,
confirmed directly, dead code in every build configuration this project currently produces
--- untestable as a consequence of that, and not a Media-specific concern.
